La deuxième relation fondamentale apparaît lorsque l’on compare l’action et la gravité.\par

Dans HMT, l’action mesure combien il se produit et conserve simultanément l’orientation de la réalisation. Il existe une section antérieure au retour et une autre postérieure. La phase peut revenir accompagnée d’une histoire de mémoire différente. Cette différence est stockée dans le rapport bidirectionnel des deux sections d’action.\par

La gravité répond de manière réciproque.\par

Le transducteur gravitationnel a la forme\par

\[
G_a(\theta)
=
\theta\frac{c_{\rm int}^{5}t_0^2}{\hbar_a}.
\]

L’action apparaît au dénominateur. Par conséquent, si, lors du changement de feuillet, \(\hbar\) est multiplié par un rapport \(R_{\rm act}\), la constante gravitationnelle est multipliée par le rapport inverse :\par

\[
\frac{\hbar_{\rm pre}}{\hbar_{\rm ret}}
=
R_{\rm act},
\qquad
\frac{G_{\rm pre}}{G_{\rm ret}}
=
R_{\rm act}^{-1}.
\]

L’action et la gravité forment ici une paire contragrédiente : lorsque l’une avance dans une direction, l’autre compense exactement ce déplacement.\par

Ce qui reste invariant est\par

\[
G_a\hbar_a.
\]

Et c’est précisément la combinaison qui détermine l’aire de Planck :\par

\[
\ell_P^2=\frac{G_a\hbar_a}{c_{\rm int}^3}.
\]

L’action mémorise le feuillet. La gravité répond à cette mémoire. L’aire reste invariante.\par

C’est une idée extraordinairement féconde. L’aire de Planck émerge comme l’objet qui subsiste lorsque l’action et la gravité se transforment réciproquement.\par

On pourrait le dire ainsi : \(\hbar\) conserve l’orientation dynamique ; \(G\) conserve la compensation géométrique ; \(\ell_P^2\) conserve ce sur quoi les deux réalisations s’accordent.\par

L’aire est donc une monnaie d’échange entre physique quantique et géométrie, parce que la transformation bidirectionnelle de \(\hbar\) et de \(G\) sélectionne exactement cette combinaison comme invariant.\par

Cette réciprocité permet de comprendre autrement la relation entre information et gravité. L’information exige une distinction entre les états ; l’action mesure le coût de la transition entre eux. Mais si cette transition doit acquérir une géométrie, la réponse gravitationnelle change de manière inverse pour maintenir stable le quantum d’aire.\par

La géométrie incorpore l’information de phase par une compensation exacte.\par

Selon une lecture machienne, c’est déjà un résultat très suggestif. Une grandeur locale de géométrie est déterminée relationnellement par l’orientation et le retour de l’état complet. L’aire invariante restreint cette dépendance relationnelle et en fixe la cohérence.\par

La version HMT du principe relationnel consiste en quelque chose de plus précis que la maxime « la matière détermine la géométrie » : l’orientation de l’action et la réponse gravitationnelle se codéterminent de telle manière que la géométrie élémentaire reste invariante.\par

L’action contient l’histoire. La gravité ajuste la réponse. L’aire conserve l’accord.\par

